% Moved verbatim from the main text during the length revision
% (was common_body.tex lines 25-54). Content unchanged.
\subsection{Post-hoc OOD Detection and Training-Time Enhancements}
\label{sec:rw-ood}

Post-hoc out-of-distribution (OOD) detection methods are typically grouped by whether they read the classifier's output layer or its internal representation. Maximum softmax probability \citep{hendrycks2017baseline} and its temperature-scaled and input-perturbed successor ODIN \citep{liang2018enhancing} score an input from the output logits directly; energy-based scoring \citep{liu2020energy} aggregates the full logit vector rather than its maximum; ReAct \citep{sun2021react} rectifies unusually high activations before scoring; ViM \citep{wang2022vim} combines a logit-space term with a residual computed in feature space; and activation-shaping \citep{djurisic2023ash} prunes a large fraction of late-layer activations at inference time without using any training-set statistics. Complementary training-time approaches can instead reshape representations or logits so that subsequent post-hoc OOD scores become more effective. Extended Logit Normalization (ELogitNorm) \citep{ding2025enhancing}, for example, introduces a feature-distance-aware training objective designed to improve a range of post-hoc OOD detectors while preserving in-distribution classification performance. Distance-based scorers treat the penultimate-layer representation as the object of interest: Mahalanobis distance \citep{lee2018simple} fits class-conditional Gaussians and scores an input by its distance to the nearest class mean; pooled $k$-nearest-neighbor distance \citep{sun2022knn} replaces this parametric fit with a non-parametric density estimate; nearest-neighbor guidance \citep{park2023nnguide} instead uses a $k$-NN-derived correction term to reduce classifier overconfidence rather than to replace the classifier score outright; and Gram-matrix scoring \citep{sastry2020gram} extracts still other structural summaries of the same representation. Benchmarking across this family, surveyed for the convolutional setting by \citet{yang2024oodsurvey} and, as the field has moved toward large pretrained encoders, for the vision-language and foundation-model setting by \citet{miyai2024generalized}, typically compares detectors across architectures and datasets under a fixed, unperturbed training regime. This study instead holds architecture and dataset fixed and asks whether the same detector's behavior on the same representation changes when the representation itself is deliberately reshaped by training -- a question orthogonal to, and not addressed by, cross-detector benchmarking.

\subsection{Domain-Adversarial Training and Representation Disentanglement}
\label{sec:rw-disentanglement}

Domain-adversarial training \citep{ganin2015unsupervised, ganin2016domain} learns a representation that a discriminator cannot use to identify the input's domain, by reversing the discriminator's gradient during backpropagation. Related approaches pursue a similar goal through different mechanisms: domain separation networks \citep{bousmalis2016domain} factor a representation into shared and domain-private subspaces and use a difference loss that imposes a soft subspace-orthogonality constraint between them; correlation-alignment methods \citep{sun2016coral} instead match second-order feature statistics across domains without an adversarial objective; and adversarial discriminative adaptation \citep{tzeng2017adversarial} applies the same discriminator-based principle asymmetrically across a source and target encoder. A related but distinct family, disentangled representation learning \citep{higgins2017betavae}, seeks to factor a representation into semantically separable generative factors rather than to suppress one nuisance factor specifically; \citet{locatello2019challenging} shows that such disentanglement is not identifiable from unlabeled data without further assumptions, a caution that does not directly bear on the supervised, single-nuisance-variable setting audited here but that motivates treating any claimed disentanglement as an empirical question rather than an assumption. All of these methods are evaluated in their originating literature primarily by downstream task accuracy or by a discriminator's inability to recover the suppressed variable; none of the above is evaluated, to our knowledge, by whether a distance-based reliability estimator built on the resulting representation continues to function as it did before the intervention, which is the question this study takes up.

\subsection{Shortcut Learning and Spurious Correlations}
\label{sec:rw-shortcut}

The failure mode motivating disentanglement -- a classifier exploiting a nuisance correlate of the label rather than task-relevant structure -- has been documented repeatedly in clinical imaging. Chest radiograph classifiers have been shown to encode hospital- and scanner-specific artefacts \citep{zech2018variable}; COVID-19 detectors trained on heterogeneous radiograph sources were found to rely on scanner-specific cues detectable by saliency attribution rather than on pathology \citep{degrave2021ai}; and in dermatology specifically, surgical skin markings have been shown to inflate a model's predicted probability of malignancy \citep{winkler2019association}, while \citet{bissoto2020debiasing} report that ISIC-trained classifiers retain above-chance accuracy even when the lesion region itself is masked, implicating acquisition-context cues directly. \citet{geirhos2020shortcut} synthesizes this literature under the term shortcut learning and argues that standard accuracy metrics provide no warning of it. This body of work establishes that shortcut reliance is a real and recurring failure mode and motivates disentanglement as a countermeasure; it does not, however, examine what such a countermeasure does to a reliability estimator layered on top of the resulting representation, which is the gap this study addresses.

\subsection{Probing Classifiers and Representation Analysis}
\label{sec:rw-probing}

Linear probing -- training a classifier on a frozen representation to measure what that representation linearly encodes -- was introduced as a diagnostic for what intermediate layers of a network represent \citep{alain2016understanding} and has since been used across modalities to characterize learned representations independently of the objective that produced them \citep{hewitt2019designing}. \citet{belinkov2022probing} surveys this literature and cautions specifically against over-interpreting a successful probe as evidence that the original model computes or uses the probed information in the way the probe accesses it. This study's use of probing follows that caution directly: probe AUROC here is offered only as a measurement of what is linearly and non-linearly recoverable from an embedding, independent of any distance-based scorer, and is not offered as a claim about the disentanglement objective's internal computation.

\subsection{Representation Geometry and Estimator Performance}
\label{sec:rw-geometry}

A separate line of work characterizes the geometry of learned representations directly, rather than what they linearly encode. Intrinsic dimensionality, estimated by nearest-neighbor methods such as the maximum-likelihood estimator of \citet{levina2004maximum}, the two-nearest-neighbor estimator of \citet{facco2017estimating}, or applied to deep network representations by \citet{ansuini2019intrinsic}, is one such summary; the conditioning of the class-covariance matrix used to fit a Mahalanobis estimator is another. Across a range of frozen vision foundation-model backbones, \citet{janiak2026geometry} report that intrinsic dimensionality and within-class spectral decay track how well Mahalanobis-style detectors perform, characterizing this relationship as an association observed across backbones rather than as a causal claim. That association is measured between representations that differ simultaneously in pretraining corpus, objective, architecture, and scale, and does not by itself establish what happens when the same geometric quantities are moved within a single, fixed representation by a specific training intervention -- the question this study asks for one such intervention, domain-adversarial disentanglement, and one such geometric summary, condition number, among others.

\subsection{Reliability Estimation in Medical Imaging Under Distribution Shift}
\label{sec:rw-medical}

That predictive uncertainty degrades under dataset shift, and that this degradation is consequential for clinical decision-making, is established broadly by \citet{ovadia2019can} and discussed specifically for clinical deployment by \citet{finlayson2021clinician}. In dermatology, out-of-distribution detection for skin lesion images has been examined with deep isolation-forest scoring \citep{li2020deepif} and benchmarked systematically across covariate-shifted, near-OOD, and far-OOD regimes \citep{gutbrod2025openmibood}; deep skin-cancer classifiers have also been shown to remain overconfident on out-of-distribution inputs \citep{pacheco2020outofdistribution}. A related and largely orthogonal source of distributional heterogeneity in this domain is skin tone: \citet{daneshjou2022disparities} curated a diverse, pathologically confirmed clinical image set and showed that dermatology AI models trained on conventional benchmarks perform substantially worse on darker skin tones and uncommon diseases, a gap this study's own dataset pair does not disentangle from acquisition-device shift (Section~\ref{sec:limitations}). The field has also begun moving from task-specific convolutional classifiers toward large pretrained encoders used as frozen feature extractors in dermatology specifically \citep{yan2025panderm}; whether the geometry-reliability dissociation reported here persists in such representations is a question this study's design does not address but that motivates direct follow-up. That literature evaluates reliability estimators against naturally occurring distribution shift; it does not evaluate what happens to those estimators when the representation they operate on is itself the target of an explicit debiasing intervention, which is the setting audited here. This study also does not evaluate calibration \citep{guo2017calibration} or selective prediction \citep{elyaniv2010foundations, geifman2017selective}, nor distribution-free alternatives to both such as conformal prediction \citep{angelopoulos2023conformal}, directly -- these are natural extensions of the reliability question addressed here and are noted as such in Section~\ref{sec:limitations}.
